\documentclass[border=12pt]{standalone}
\usepackage[utf8]{inputenc}
\usepackage[T1]{fontenc}
\usepackage[american, siunitx]{circuitikz}

\sisetup{output-decimal-marker={,}}

\ctikzset{
    resistors/scale=0.8,
    diodes/scale=0.8
}

\begin{document}
\begin{circuitikz}[font=\sffamily\small]

    % =========================================================================
    % 1. BLOCO ESQUERDO: ESP32 DEVKITC V4 (U1)
    % =========================================================================
    \draw[thick, fill=gray!5] (-5.2, -2.6) rectangle (0, 2.6);
    \node[align=center, font=\sffamily\bfseries] at (-3.0, 0.3) {U1\\ESP32};
    \node[font=\sffamily\scriptsize, text=gray!70!black] at (-3.0, -0.6) {(DevKitC V4)};

    % Alimentacao de U1 (+3,3V e GND)
    \draw (-2.6, 2.6) -- (-2.6, 3.4) node[vcc]{+3,3V};
    \node[below, font=\ttfamily\footnotesize] at (-2.6, 2.6) {3V3};

    \draw (-2.6, -2.6) -- (-2.6, -3.4) node[ground]{};
    \node[above, font=\ttfamily\footnotesize] at (-2.6, -2.6) {GND};

    % =========================================================================
    % 2. BLOCO DIREITO: RESISTORES (220R) E LED RGB CATODO COMUM (LED1)
    % =========================================================================
    % Encapsulamento do LED RGB
    \draw[dashed, thick, fill=gray!5] (4.4, -2.3) rectangle (8.4, 2.3);
    \node[above, font=\sffamily\bfseries\footnotesize] at (6.4, 2.35) {LED1 -- RGB (Cátodo Comum)};

    % Conexoes GPIO -> Resistores -> Anodos (R, G, B)
    \foreach \y/\pin/\rname/\cname in {
        1.5/GPIO5/R1/R,
        0.0/GPIO17/R2/G,
        -1.5/GPIO16/R3/B%
    }{
        \node[left, font=\ttfamily\scriptsize] at (0, \y) {\pin};
        \draw (0, \y) -- (0.8, \y)
            to[R, l^=\rname, a_=220<\ohm>] (3.6, \y)
            -- (4.8, \y)
            to[leDo, l^=\cname] (7.2, \y)
            -- (7.8, \y);
    }

    % Barramento comum de catodo (COM -> GND, sem ponto no canto superior)
    \draw (7.8, 1.5) -- (7.8, -3.2) node[ground]{};
    \node[circ] at (7.8, 0.0) {};
    \node[circ] at (7.8, -1.5) {};
    \node[right, font=\ttfamily\scriptsize] at (7.8, -2.65) {COM};

\end{circuitikz}
\end{document}